% --- Archon physics formalization source begin ---
% archon:physics
% archon:covers PhyXMiniProblems/problem_phyx_mini_0907.lean
% archon:source-report reports/phyx_mini/problem_phyx_mini_0907.source.json

\chapter{Physics problem phyx\_mini\_0907}
\label{ch:PhyXMiniProblems_problem_phyx_mini_0907}

\paragraph{Problem source.}
\#\# Physical scenario

The force on the -1.0 nC charge is as shown in Figure

\#\# Question

What is the magnitude of this force?

\#\# Answer choices

- A: A: 7.25 \textbackslash{}times 10\textasciicircum{}\{-4\}N

- B: B: 6.35 \textbackslash{}times 10\textasciicircum{}\{-4\}N

- C: C: 1.7 \textbackslash{}times 10\textasciicircum{}\{-4\}N

- D: D: 1.33 \textbackslash{}times 10\textasciicircum{}\{-4\}N

\#\# Figure caption (auxiliary; use the image as primary evidence)

The image is a diagram illustrating an arrangement of three charged objects. The objects and their relationships are described as follows:

1. **Objects:**
   - A red circle labeled with a "+", representing a positive charge of \textbackslash{}(10 \textbackslash{}, \textbackslash{}text\{nC\}\textbackslash{}).
   - A cyan circle labeled with a "−", representing a negative charge of \textbackslash{}(-1.0 \textbackslash{}, \textbackslash{}text\{nC\}\textbackslash{}).
   - A gray circle labeled with "q", the value of which is unspecified.

2. **Arrangement and Measurements:**
   - The three charges form a triangle.
   - The distance between the gray circle "q" and the red positive charge is labeled as \textbackslash{}(5.0 \textbackslash{}, \textbackslash{}text\{cm\}\textbackslash{}).
   - The angle at the red positive charge is \textbackslash{}(60\textasciicircum{}\textbackslash{}circ\textbackslash{}).
   - The angle at the gray circle "q" is \textbackslash{}(30\textasciicircum{}\textbackslash{}circ\textbackslash{}).

3. **Forces and Vectors:**
   - There is a downward-pointing red arrow labeled \textbackslash{}(\textbackslash{}vec\{F\}\textbackslash{}) starting from the negative charge, likely indicating a force vector.

4. **Angles and Geometry:**
   - A right angle is indicated where the perpendicular line drops down from the negative charge to the base of the triangle, formed by the line segment connecting the positive charge and the gray circle \textbackslash{}(q\textbackslash{}).

The diagram emphasizes the geometric and electrostatic relationships among the three charges.

\#\# Dataset metadata

Electrostatics; Physical Model Grounding Reasoning, Spatial Relation Reasoning

\paragraph{Recorded answer/context.}
C: C: 1.7 \textbackslash{}times 10\textasciicircum{}\{-4\}N

\paragraph{Figure/image path.}
/root/proposal\_for\_physic/science-mango/phyx\_mini\_run/phyx\_data/test\_image/907.png

\paragraph{Formalization target.}
create a compiling Lean file with sorry bodies at `PhyXMiniProblems/problem\_phyx\_mini\_0907.lean`. The Lean declarations must preserve the physical quantities, dimensions or dimensional roles, figure labels, governing-law hypotheses, and final relation expressed by this problem.
Use Mathlib/Physlib names found through LeanExplore where available. If a physics API is missing, introduce faithful local abstractions rather than scalar placeholder aliases.

\begin{theorem}[Physics formalization target]
\label{thm:physics:phyx_mini_0907:target}
\lean{PhyXMiniProblems.ProblemPhyXMini0907.problem_phyx_mini_0907}
\uses{def:physics:phyx-mini-0907:phyxminiproblems-problemphyxmini0907-forcemagnitudeinnewtons, def:physics:phyx-mini-0907:phyxminiproblems-problemphyxmini0907-threechargetrianglesetup, def:physics:phyx-mini-0907:phyxminiproblems-problemphyxmini0907-matchesthreepointchargescenario, def:physics:phyx-mini-0907:phyxminiproblems-problemphyxmini0907-matchessuppliedthreechargetrianglefigure, def:physics:phyx-mini-0907:phyxminiproblems-problemphyxmini0907-matchesproblemphysicalreadouts, def:physics:phyx-mini-0907:phyxminiproblems-problemphyxmini0907-satisfiesdisplayedtrianglegeometry, def:physics:phyx-mini-0907:phyxminiproblems-problemphyxmini0907-usesschoolcoulombconstant, def:physics:phyx-mini-0907:phyxminiproblems-problemphyxmini0907-hasphysicalelectrostaticparameters, def:physics:phyx-mini-0907:phyxminiproblems-problemphyxmini0907-calibratesdisplayeddownwardforcedirection, def:physics:phyx-mini-0907:phyxminiproblems-problemphyxmini0907-satisfiesvectorcoulombforcelaw, def:physics:phyx-mini-0907:phyxminiproblems-problemphyxmini0907-satisfieselectrostaticforcesuperposition, def:physics:phyx-mini-0907:phyxminiproblems-problemphyxmini0907-satisfiesresultantforcemagnitudelaw, def:physics:phyx-mini-0907:phyxminiproblems-problemphyxmini0907-displayedforcemagnitudeinnewtons, def:physics:phyx-mini-0907:phyxminiproblems-problemphyxmini0907-recordeddatasetanswer, def:physics:phyx-mini-0907:phyxminiproblems-problemphyxmini0907-roundstoatresolution, def:physics:phyx-mini-0907:phyxminiproblems-problemphyxmini0907-isuniqueclosestdisplayedforcechoice}
The assigned autoformalize agent should translate this physics problem into Lean declarations in the covered file, with theorem and lemma proof bodies written as `by sorry`.
\end{theorem}
\begin{proof}
This is an autoformalization task, not a proof task. Produce statements that can later be proved by the physics prover without weakening the physical model.
\end{proof}

% --- Archon declaration topology begin ---
\section{Declaration topology}
\begin{definition}[force Dimension]
  \label{def:physics:phyx-mini-0907:phyxminiproblems-problemphyxmini0907-forcedimension}
  \lean{PhyXMiniProblems.ProblemPhyXMini0907.forceDimension}
  The physical dimension M L T⁻² of force
\end{definition}
\begin{proof}
  This is the declared model definition of force Dimension.
\end{proof}
\begin{definition}[Length Quantity]
  \label{def:physics:phyx-mini-0907:phyxminiproblems-problemphyxmini0907-lengthquantity}
  \lean{PhyXMiniProblems.ProblemPhyXMini0907.LengthQuantity}
  A nonnegative, unit-independent physical length
\end{definition}
\begin{proof}
  This is the declared model definition of Length Quantity.
\end{proof}
\begin{definition}[Signed Charge Quantity]
  \label{def:physics:phyx-mini-0907:phyxminiproblems-problemphyxmini0907-signedchargequantity}
  \lean{PhyXMiniProblems.ProblemPhyXMini0907.SignedChargeQuantity}
  A signed, unit-independent electric charge
\end{definition}
\begin{proof}
  This is the declared model definition of Signed Charge Quantity.
\end{proof}
\begin{definition}[Planar Position Quantity]
  \label{def:physics:phyx-mini-0907:phyxminiproblems-problemphyxmini0907-planarpositionquantity}
  \lean{PhyXMiniProblems.ProblemPhyXMini0907.PlanarPositionQuantity}
  A unit-independent position vector in the plane of the figure
\end{definition}
\begin{proof}
  This is the declared model definition of Planar Position Quantity.
\end{proof}
\begin{definition}[Planar Force Quantity]
  \label{def:physics:phyx-mini-0907:phyxminiproblems-problemphyxmini0907-planarforcequantity}
  \lean{PhyXMiniProblems.ProblemPhyXMini0907.PlanarForceQuantity}
  \uses{def:physics:phyx-mini-0907:phyxminiproblems-problemphyxmini0907-forcedimension}
  A unit-independent force vector in the plane of the figure
\end{definition}
\begin{proof}
  This is the declared model definition of Planar Force Quantity.
\end{proof}
\begin{definition}[Force Magnitude Quantity]
  \label{def:physics:phyx-mini-0907:phyxminiproblems-problemphyxmini0907-forcemagnitudequantity}
  \lean{PhyXMiniProblems.ProblemPhyXMini0907.ForceMagnitudeQuantity}
  \uses{def:physics:phyx-mini-0907:phyxminiproblems-problemphyxmini0907-forcedimension}
  A nonnegative, unit-independent physical force magnitude
\end{definition}
\begin{proof}
  This is the declared model definition of Force Magnitude Quantity.
\end{proof}
\begin{definition}[x Axis]
  \label{def:physics:phyx-mini-0907:phyxminiproblems-problemphyxmini0907-xaxis}
  \lean{PhyXMiniProblems.ProblemPhyXMini0907.xAxis}
  Coordinate 0, pointing right along the horizontal base
\end{definition}
\begin{proof}
  This is the declared model definition of x Axis.
\end{proof}
\begin{definition}[y Axis]
  \label{def:physics:phyx-mini-0907:phyxminiproblems-problemphyxmini0907-yaxis}
  \lean{PhyXMiniProblems.ProblemPhyXMini0907.yAxis}
  Coordinate 1, pointing upward in the image
\end{definition}
\begin{proof}
  This is the declared model definition of y Axis.
\end{proof}
\begin{definition}[length In Meters]
  \label{def:physics:phyx-mini-0907:phyxminiproblems-problemphyxmini0907-lengthinmeters}
  \lean{PhyXMiniProblems.ProblemPhyXMini0907.lengthInMeters}
  \uses{def:physics:phyx-mini-0907:phyxminiproblems-problemphyxmini0907-lengthquantity}
  Coherent-SI readout of a physical length in metres
\end{definition}
\begin{proof}
  This is the declared model definition of length In Meters.
\end{proof}
\begin{definition}[length In Centimeters]
  \label{def:physics:phyx-mini-0907:phyxminiproblems-problemphyxmini0907-lengthincentimeters}
  \lean{PhyXMiniProblems.ProblemPhyXMini0907.lengthInCentimeters}
  \uses{def:physics:phyx-mini-0907:phyxminiproblems-problemphyxmini0907-lengthquantity, def:physics:phyx-mini-0907:phyxminiproblems-problemphyxmini0907-lengthinmeters}
  Centimetre readout used by the base label in the image
\end{definition}
\begin{proof}
  This is the declared model definition of length In Centimeters.
\end{proof}
\begin{definition}[charge In Coulombs]
  \label{def:physics:phyx-mini-0907:phyxminiproblems-problemphyxmini0907-chargeincoulombs}
  \lean{PhyXMiniProblems.ProblemPhyXMini0907.chargeInCoulombs}
  \uses{def:physics:phyx-mini-0907:phyxminiproblems-problemphyxmini0907-signedchargequantity}
  Coherent-SI readout of a signed charge in coulombs
\end{definition}
\begin{proof}
  This is the declared model definition of charge In Coulombs.
\end{proof}
\begin{definition}[charge In Nanocoulombs]
  \label{def:physics:phyx-mini-0907:phyxminiproblems-problemphyxmini0907-chargeinnanocoulombs}
  \lean{PhyXMiniProblems.ProblemPhyXMini0907.chargeInNanocoulombs}
  \uses{def:physics:phyx-mini-0907:phyxminiproblems-problemphyxmini0907-signedchargequantity, def:physics:phyx-mini-0907:phyxminiproblems-problemphyxmini0907-chargeincoulombs}
  Nanocoulomb readout used by the two numerical charge labels
\end{definition}
\begin{proof}
  This is the declared model definition of charge In Nanocoulombs.
\end{proof}
\begin{definition}[position Vector In Meters]
  \label{def:physics:phyx-mini-0907:phyxminiproblems-problemphyxmini0907-positionvectorinmeters}
  \lean{PhyXMiniProblems.ProblemPhyXMini0907.positionVectorInMeters}
  \uses{def:physics:phyx-mini-0907:phyxminiproblems-problemphyxmini0907-planarpositionquantity}
  Coherent-SI Cartesian position vector, with coordinates in metres
\end{definition}
\begin{proof}
  This is the declared model definition of position Vector In Meters.
\end{proof}
\begin{definition}[force Vector In Newtons]
  \label{def:physics:phyx-mini-0907:phyxminiproblems-problemphyxmini0907-forcevectorinnewtons}
  \lean{PhyXMiniProblems.ProblemPhyXMini0907.forceVectorInNewtons}
  \uses{def:physics:phyx-mini-0907:phyxminiproblems-problemphyxmini0907-planarforcequantity}
  Coherent-SI Cartesian force vector, with components in newtons
\end{definition}
\begin{proof}
  This is the declared model definition of force Vector In Newtons.
\end{proof}
\begin{definition}[force Magnitude In Newtons]
  \label{def:physics:phyx-mini-0907:phyxminiproblems-problemphyxmini0907-forcemagnitudeinnewtons}
  \lean{PhyXMiniProblems.ProblemPhyXMini0907.forceMagnitudeInNewtons}
  \uses{def:physics:phyx-mini-0907:phyxminiproblems-problemphyxmini0907-forcemagnitudequantity}
  Coherent-SI readout of a force magnitude in newtons
\end{definition}
\begin{proof}
  This is the declared model definition of force Magnitude In Newtons.
\end{proof}
\begin{definition}[angle In Degrees]
  \label{def:physics:phyx-mini-0907:phyxminiproblems-problemphyxmini0907-angleindegrees}
  \lean{PhyXMiniProblems.ProblemPhyXMini0907.angleInDegrees}
  Read the principal representative of an angle in degrees
\end{definition}
\begin{proof}
  This is the declared model definition of angle In Degrees.
\end{proof}
\begin{definition}[Particle]
  \label{def:physics:phyx-mini-0907:phyxminiproblems-problemphyxmini0907-particle}
  \lean{PhyXMiniProblems.ProblemPhyXMini0907.Particle}
  The three charged objects, named by their roles in the supplied image
\end{definition}
\begin{proof}
  This is the declared model definition of Particle.
\end{proof}
\begin{definition}[Force Source]
  \label{def:physics:phyx-mini-0907:phyxminiproblems-problemphyxmini0907-forcesource}
  \lean{PhyXMiniProblems.ProblemPhyXMini0907.ForceSource}
  The two particles exerting electrostatic force on the target charge
\end{definition}
\begin{proof}
  This is the declared model definition of Force Source.
\end{proof}
\begin{definition}[to Particle]
  \label{def:physics:phyx-mini-0907:phyxminiproblems-problemphyxmini0907-forcesource-toparticle}
  \lean{PhyXMiniProblems.ProblemPhyXMini0907.ForceSource.toParticle}
  \uses{def:physics:phyx-mini-0907:phyxminiproblems-problemphyxmini0907-particle, def:physics:phyx-mini-0907:phyxminiproblems-problemphyxmini0907-forcesource}
  Regard a force source as one of the three pictured particles
\end{definition}
\begin{proof}
  This is the declared model definition of to Particle.
\end{proof}
\begin{definition}[Charge Site]
  \label{def:physics:phyx-mini-0907:phyxminiproblems-problemphyxmini0907-chargesite}
  \lean{PhyXMiniProblems.ProblemPhyXMini0907.ChargeSite}
  Geometric site occupied by each charge in the triangle
\end{definition}
\begin{proof}
  This is the declared model definition of Charge Site.
\end{proof}
\begin{definition}[expected Particle Site]
  \label{def:physics:phyx-mini-0907:phyxminiproblems-problemphyxmini0907-expectedparticlesite}
  \lean{PhyXMiniProblems.ProblemPhyXMini0907.expectedParticleSite}
  \uses{def:physics:phyx-mini-0907:phyxminiproblems-problemphyxmini0907-particle, def:physics:phyx-mini-0907:phyxminiproblems-problemphyxmini0907-chargesite}
  Site at which each named particle appears
\end{definition}
\begin{proof}
  This is the declared model definition of expected Particle Site.
\end{proof}
\begin{definition}[expected Particle Label]
  \label{def:physics:phyx-mini-0907:phyxminiproblems-problemphyxmini0907-expectedparticlelabel}
  \lean{PhyXMiniProblems.ProblemPhyXMini0907.expectedParticleLabel}
  \uses{def:physics:phyx-mini-0907:phyxminiproblems-problemphyxmini0907-particle}
  Text printed beside each charge circle
\end{definition}
\begin{proof}
  This is the declared model definition of expected Particle Label.
\end{proof}
\begin{definition}[expected Printed Charge Nanocoulombs]
  \label{def:physics:phyx-mini-0907:phyxminiproblems-problemphyxmini0907-expectedprintedchargenanocoulombs}
  \lean{PhyXMiniProblems.ProblemPhyXMini0907.expectedPrintedChargeNanocoulombs}
  \uses{def:physics:phyx-mini-0907:phyxminiproblems-problemphyxmini0907-particle}
  Numerical nanocoulomb data printed in the image; q is unspecified
\end{definition}
\begin{proof}
  This is the declared model definition of expected Printed Charge Nanocoulombs.
\end{proof}
\begin{definition}[Figure Charge Sign]
  \label{def:physics:phyx-mini-0907:phyxminiproblems-problemphyxmini0907-figurechargesign}
  \lean{PhyXMiniProblems.ProblemPhyXMini0907.FigureChargeSign}
  Sign glyph drawn inside a charge circle
\end{definition}
\begin{proof}
  This is the declared model definition of Figure Charge Sign.
\end{proof}
\begin{definition}[expected Charge Sign]
  \label{def:physics:phyx-mini-0907:phyxminiproblems-problemphyxmini0907-expectedchargesign}
  \lean{PhyXMiniProblems.ProblemPhyXMini0907.expectedChargeSign}
  \uses{def:physics:phyx-mini-0907:phyxminiproblems-problemphyxmini0907-particle, def:physics:phyx-mini-0907:phyxminiproblems-problemphyxmini0907-figurechargesign}
  The unknown q circle has no sign glyph in the supplied image
\end{definition}
\begin{proof}
  This is the declared model definition of expected Charge Sign.
\end{proof}
\begin{definition}[Figure Charge Color]
  \label{def:physics:phyx-mini-0907:phyxminiproblems-problemphyxmini0907-figurechargecolor}
  \lean{PhyXMiniProblems.ProblemPhyXMini0907.FigureChargeColor}
  Fill colour of a charge circle in the supplied raster
\end{definition}
\begin{proof}
  This is the declared model definition of Figure Charge Color.
\end{proof}
\begin{definition}[expected Charge Color]
  \label{def:physics:phyx-mini-0907:phyxminiproblems-problemphyxmini0907-expectedchargecolor}
  \lean{PhyXMiniProblems.ProblemPhyXMini0907.expectedChargeColor}
  \uses{def:physics:phyx-mini-0907:phyxminiproblems-problemphyxmini0907-particle, def:physics:phyx-mini-0907:phyxminiproblems-problemphyxmini0907-figurechargecolor}
  Expected circle colour of each particle
\end{definition}
\begin{proof}
  This is the declared model definition of expected Charge Color.
\end{proof}
\begin{definition}[Base Vertex]
  \label{def:physics:phyx-mini-0907:phyxminiproblems-problemphyxmini0907-basevertex}
  \lean{PhyXMiniProblems.ProblemPhyXMini0907.BaseVertex}
  The two printed acute angles at the endpoints of the base
\end{definition}
\begin{proof}
  This is the declared model definition of Base Vertex.
\end{proof}
\begin{definition}[expected Base Angle Degrees]
  \label{def:physics:phyx-mini-0907:phyxminiproblems-problemphyxmini0907-expectedbaseangledegrees}
  \lean{PhyXMiniProblems.ProblemPhyXMini0907.expectedBaseAngleDegrees}
  \uses{def:physics:phyx-mini-0907:phyxminiproblems-problemphyxmini0907-basevertex}
  Angle in degrees printed at each base vertex
\end{definition}
\begin{proof}
  This is the declared model definition of expected Base Angle Degrees.
\end{proof}
\begin{definition}[Triangle Side]
  \label{def:physics:phyx-mini-0907:phyxminiproblems-problemphyxmini0907-triangleside}
  \lean{PhyXMiniProblems.ProblemPhyXMini0907.TriangleSide}
  The three dashed sides forming the charge triangle
\end{definition}
\begin{proof}
  This is the declared model definition of Triangle Side.
\end{proof}
\begin{definition}[Arrow Direction]
  \label{def:physics:phyx-mini-0907:phyxminiproblems-problemphyxmini0907-arrowdirection}
  \lean{PhyXMiniProblems.ProblemPhyXMini0907.ArrowDirection}
  Direction of the single red force arrow
\end{definition}
\begin{proof}
  This is the declared model definition of Arrow Direction.
\end{proof}
\begin{definition}[Three Charge Triangle Figure]
  \label{def:physics:phyx-mini-0907:phyxminiproblems-problemphyxmini0907-threechargetrianglefigure}
  \lean{PhyXMiniProblems.ProblemPhyXMini0907.ThreeChargeTriangleFigure}
  \uses{def:physics:phyx-mini-0907:phyxminiproblems-problemphyxmini0907-particle, def:physics:phyx-mini-0907:phyxminiproblems-problemphyxmini0907-chargesite, def:physics:phyx-mini-0907:phyxminiproblems-problemphyxmini0907-figurechargesign, def:physics:phyx-mini-0907:phyxminiproblems-problemphyxmini0907-figurechargecolor, def:physics:phyx-mini-0907:phyxminiproblems-problemphyxmini0907-basevertex, def:physics:phyx-mini-0907:phyxminiproblems-problemphyxmini0907-triangleside, def:physics:phyx-mini-0907:phyxminiproblems-problemphyxmini0907-arrowdirection}
  Literal presentation data transcribed from primary image 907.png
\end{definition}
\begin{proof}
  This is the declared model definition of Three Charge Triangle Figure.
\end{proof}
\begin{definition}[Three Charge Triangle Setup]
  \label{def:physics:phyx-mini-0907:phyxminiproblems-problemphyxmini0907-threechargetrianglesetup}
  \lean{PhyXMiniProblems.ProblemPhyXMini0907.ThreeChargeTriangleSetup}
  \uses{def:physics:phyx-mini-0907:phyxminiproblems-problemphyxmini0907-lengthquantity, def:physics:phyx-mini-0907:phyxminiproblems-problemphyxmini0907-signedchargequantity, def:physics:phyx-mini-0907:phyxminiproblems-problemphyxmini0907-planarpositionquantity, def:physics:phyx-mini-0907:phyxminiproblems-problemphyxmini0907-planarforcequantity, def:physics:phyx-mini-0907:phyxminiproblems-problemphyxmini0907-forcemagnitudequantity, def:physics:phyx-mini-0907:phyxminiproblems-problemphyxmini0907-particle, def:physics:phyx-mini-0907:phyxminiproblems-problemphyxmini0907-forcesource, def:physics:phyx-mini-0907:phyxminiproblems-problemphyxmini0907-basevertex, def:physics:phyx-mini-0907:phyxminiproblems-problemphyxmini0907-threechargetrianglefigure}
  Independent physical quantities and force observables. In particular, the unknown charge, pairwise forces, net force, and net-force magnitude are not defined from the recorded answer
\end{definition}
\begin{proof}
  This is the declared model definition of Three Charge Triangle Setup.
\end{proof}
\begin{definition}[displacement From Source To Target In Meters]
  \label{def:physics:phyx-mini-0907:phyxminiproblems-problemphyxmini0907-displacementfromsourcetotargetinmeters}
  \lean{PhyXMiniProblems.ProblemPhyXMini0907.displacementFromSourceToTargetInMeters}
  \uses{def:physics:phyx-mini-0907:phyxminiproblems-problemphyxmini0907-positionvectorinmeters, def:physics:phyx-mini-0907:phyxminiproblems-problemphyxmini0907-forcesource, def:physics:phyx-mini-0907:phyxminiproblems-problemphyxmini0907-threechargetrianglesetup}
  Displacement from a source particle to the target particle, in metres
\end{definition}
\begin{proof}
  This is the declared model definition of displacement From Source To Target In Meters.
\end{proof}
\begin{definition}[Matches Three Point Charge Scenario]
  \label{def:physics:phyx-mini-0907:phyxminiproblems-problemphyxmini0907-matchesthreepointchargescenario}
  \lean{PhyXMiniProblems.ProblemPhyXMini0907.MatchesThreePointChargeScenario}
  \uses{def:physics:phyx-mini-0907:phyxminiproblems-problemphyxmini0907-threechargetrianglesetup}
  All three small charged objects are idealized as point charges
\end{definition}
\begin{proof}
  This is the declared model definition of Matches Three Point Charge Scenario.
\end{proof}
\begin{definition}[Matches Supplied Three Charge Triangle Figure]
  \label{def:physics:phyx-mini-0907:phyxminiproblems-problemphyxmini0907-matchessuppliedthreechargetrianglefigure}
  \lean{PhyXMiniProblems.ProblemPhyXMini0907.MatchesSuppliedThreeChargeTriangleFigure}
  \uses{def:physics:phyx-mini-0907:phyxminiproblems-problemphyxmini0907-expectedparticlesite, def:physics:phyx-mini-0907:phyxminiproblems-problemphyxmini0907-expectedparticlelabel, def:physics:phyx-mini-0907:phyxminiproblems-problemphyxmini0907-expectedprintedchargenanocoulombs, def:physics:phyx-mini-0907:phyxminiproblems-problemphyxmini0907-expectedchargesign, def:physics:phyx-mini-0907:phyxminiproblems-problemphyxmini0907-expectedchargecolor, def:physics:phyx-mini-0907:phyxminiproblems-problemphyxmini0907-expectedbaseangledegrees, def:physics:phyx-mini-0907:phyxminiproblems-problemphyxmini0907-threechargetrianglesetup}
  Literal qualitative content of the primary raster. The arrow information records only which force is depicted and its direction, never its magnitude
\end{definition}
\begin{proof}
  This is the declared model definition of Matches Supplied Three Charge Triangle Figure.
\end{proof}
\begin{definition}[Matches Problem Physical Readouts]
  \label{def:physics:phyx-mini-0907:phyxminiproblems-problemphyxmini0907-matchesproblemphysicalreadouts}
  \lean{PhyXMiniProblems.ProblemPhyXMini0907.MatchesProblemPhysicalReadouts}
  \uses{def:physics:phyx-mini-0907:phyxminiproblems-problemphyxmini0907-lengthincentimeters, def:physics:phyx-mini-0907:phyxminiproblems-problemphyxmini0907-chargeinnanocoulombs, def:physics:phyx-mini-0907:phyxminiproblems-problemphyxmini0907-angleindegrees, def:physics:phyx-mini-0907:phyxminiproblems-problemphyxmini0907-threechargetrianglesetup}
  Calibration of the numerical figure labels to physical quantities. The unknown charge q deliberately has no numerical calibration here
\end{definition}
\begin{proof}
  This is the declared model definition of Matches Problem Physical Readouts.
\end{proof}
\begin{definition}[Satisfies Displayed Triangle Geometry]
  \label{def:physics:phyx-mini-0907:phyxminiproblems-problemphyxmini0907-satisfiesdisplayedtrianglegeometry}
  \lean{PhyXMiniProblems.ProblemPhyXMini0907.SatisfiesDisplayedTriangleGeometry}
  \uses{def:physics:phyx-mini-0907:phyxminiproblems-problemphyxmini0907-lengthinmeters, def:physics:phyx-mini-0907:phyxminiproblems-problemphyxmini0907-positionvectorinmeters, def:physics:phyx-mini-0907:phyxminiproblems-problemphyxmini0907-threechargetrianglesetup}
  A coordinate realization of the displayed 30°--60°--90° triangle. The left base charge is the origin, the right base charge is on the positive horizontal axis, and the altitude foot is one quarter of the base length from the left endpoint
\end{definition}
\begin{proof}
  This is the declared model definition of Satisfies Displayed Triangle Geometry.
\end{proof}
\begin{definition}[Uses School Coulomb Constant]
  \label{def:physics:phyx-mini-0907:phyxminiproblems-problemphyxmini0907-usesschoolcoulombconstant}
  \lean{PhyXMiniProblems.ProblemPhyXMini0907.UsesSchoolCoulombConstant}
  \uses{def:physics:phyx-mini-0907:phyxminiproblems-problemphyxmini0907-threechargetrianglesetup}
  The rounded textbook value of Coulomb's constant. Physlib's scalar is read coherently in N m²/C² in the force law below
\end{definition}
\begin{proof}
  This is the declared model definition of Uses School Coulomb Constant.
\end{proof}
\begin{definition}[Has Physical Electrostatic Parameters]
  \label{def:physics:phyx-mini-0907:phyxminiproblems-problemphyxmini0907-hasphysicalelectrostaticparameters}
  \lean{PhyXMiniProblems.ProblemPhyXMini0907.HasPhysicalElectrostaticParameters}
  \uses{def:physics:phyx-mini-0907:phyxminiproblems-problemphyxmini0907-chargeincoulombs, def:physics:phyx-mini-0907:phyxminiproblems-problemphyxmini0907-threechargetrianglesetup, def:physics:phyx-mini-0907:phyxminiproblems-problemphyxmini0907-displacementfromsourcetotargetinmeters}
  Nondegeneracy and positivity conditions selecting the physical branch
\end{definition}
\begin{proof}
  This is the declared model definition of Has Physical Electrostatic Parameters.
\end{proof}
\begin{definition}[Calibrates Displayed Downward Force Direction]
  \label{def:physics:phyx-mini-0907:phyxminiproblems-problemphyxmini0907-calibratesdisplayeddownwardforcedirection}
  \lean{PhyXMiniProblems.ProblemPhyXMini0907.CalibratesDisplayedDownwardForceDirection}
  \uses{def:physics:phyx-mini-0907:phyxminiproblems-problemphyxmini0907-xaxis, def:physics:phyx-mini-0907:phyxminiproblems-problemphyxmini0907-yaxis, def:physics:phyx-mini-0907:phyxminiproblems-problemphyxmini0907-forcevectorinnewtons, def:physics:phyx-mini-0907:phyxminiproblems-problemphyxmini0907-threechargetrianglesetup}
  The figure's downward arrow is calibrated to the independently modelled net force: its horizontal component vanishes and its vertical component is negative. This supplies direction only, not the requested magnitude
\end{definition}
\begin{proof}
  This is the declared model definition of Calibrates Displayed Downward Force Direction.
\end{proof}
\begin{definition}[Satisfies Vector Coulomb Force Law]
  \label{def:physics:phyx-mini-0907:phyxminiproblems-problemphyxmini0907-satisfiesvectorcoulombforcelaw}
  \lean{PhyXMiniProblems.ProblemPhyXMini0907.SatisfiesVectorCoulombForceLaw}
  \uses{def:physics:phyx-mini-0907:phyxminiproblems-problemphyxmini0907-chargeincoulombs, def:physics:phyx-mini-0907:phyxminiproblems-problemphyxmini0907-forcevectorinnewtons, def:physics:phyx-mini-0907:phyxminiproblems-problemphyxmini0907-threechargetrianglesetup, def:physics:phyx-mini-0907:phyxminiproblems-problemphyxmini0907-displacementfromsourcetotargetinmeters}
  For source charge qₛ, target charge qₜ, and displacement r from source to target, the force on the target is F = k qₜ qₛ r / ‖r‖³. The signed product of the charges supplies attraction or repulsion. This is a general pairwise law and contains no requested net-force magnitude
\end{definition}
\begin{proof}
  This is the declared model definition of Satisfies Vector Coulomb Force Law.
\end{proof}
\begin{definition}[Satisfies Electrostatic Force Superposition]
  \label{def:physics:phyx-mini-0907:phyxminiproblems-problemphyxmini0907-satisfieselectrostaticforcesuperposition}
  \lean{PhyXMiniProblems.ProblemPhyXMini0907.SatisfiesElectrostaticForceSuperposition}
  \uses{def:physics:phyx-mini-0907:phyxminiproblems-problemphyxmini0907-forcevectorinnewtons, def:physics:phyx-mini-0907:phyxminiproblems-problemphyxmini0907-forcesource, def:physics:phyx-mini-0907:phyxminiproblems-problemphyxmini0907-threechargetrianglesetup}
  The net force is the vector sum of the two pairwise Coulomb forces
\end{definition}
\begin{proof}
  This is the declared model definition of Satisfies Electrostatic Force Superposition.
\end{proof}
\begin{definition}[Satisfies Resultant Force Magnitude Law]
  \label{def:physics:phyx-mini-0907:phyxminiproblems-problemphyxmini0907-satisfiesresultantforcemagnitudelaw}
  \lean{PhyXMiniProblems.ProblemPhyXMini0907.SatisfiesResultantForceMagnitudeLaw}
  \uses{def:physics:phyx-mini-0907:phyxminiproblems-problemphyxmini0907-forcevectorinnewtons, def:physics:phyx-mini-0907:phyxminiproblems-problemphyxmini0907-forcemagnitudeinnewtons, def:physics:phyx-mini-0907:phyxminiproblems-problemphyxmini0907-threechargetrianglesetup}
  The physical magnitude observable is the norm of the net force vector
\end{definition}
\begin{proof}
  This is the declared model definition of Satisfies Resultant Force Magnitude Law.
\end{proof}
\begin{lemma}[source target distances from triangle]
  \label{lem:physics:phyx-mini-0907:phyxminiproblems-problemphyxmini0907-source-target-distances-from-triangle}
  \lean{PhyXMiniProblems.ProblemPhyXMini0907.source_target_distances_from_triangle}
  \uses{def:physics:phyx-mini-0907:phyxminiproblems-problemphyxmini0907-lengthinmeters, def:physics:phyx-mini-0907:phyxminiproblems-problemphyxmini0907-threechargetrianglesetup, def:physics:phyx-mini-0907:phyxminiproblems-problemphyxmini0907-displacementfromsourcetotargetinmeters, def:physics:phyx-mini-0907:phyxminiproblems-problemphyxmini0907-satisfiesdisplayedtrianglegeometry}
  The coordinate realization gives the two source-to-target distances of the 30°--60°--90° triangle
\end{lemma}
\begin{proof}
  The conclusion follows from the governing relations and figure readouts named in the statement of source target distances from triangle.
\end{proof}
\begin{definition}[Answer Choice]
  \label{def:physics:phyx-mini-0907:phyxminiproblems-problemphyxmini0907-answerchoice}
  \lean{PhyXMiniProblems.ProblemPhyXMini0907.AnswerChoice}
  Labels attached to the four force-magnitude choices in the source
\end{definition}
\begin{proof}
  This is the declared model definition of Answer Choice.
\end{proof}
\begin{definition}[displayed Force Magnitude In Newtons]
  \label{def:physics:phyx-mini-0907:phyxminiproblems-problemphyxmini0907-displayedforcemagnitudeinnewtons}
  \lean{PhyXMiniProblems.ProblemPhyXMini0907.displayedForceMagnitudeInNewtons}
  \uses{def:physics:phyx-mini-0907:phyxminiproblems-problemphyxmini0907-answerchoice}
  Magnitude in newtons printed beside each displayed answer label
\end{definition}
\begin{proof}
  This is the declared model definition of displayed Force Magnitude In Newtons.
\end{proof}
\begin{definition}[recorded Dataset Answer]
  \label{def:physics:phyx-mini-0907:phyxminiproblems-problemphyxmini0907-recordeddatasetanswer}
  \lean{PhyXMiniProblems.ProblemPhyXMini0907.recordedDatasetAnswer}
  \uses{def:physics:phyx-mini-0907:phyxminiproblems-problemphyxmini0907-answerchoice}
  The answer label recorded by the supplied dataset metadata
\end{definition}
\begin{proof}
  This is the declared model definition of recorded Dataset Answer.
\end{proof}
\begin{definition}[Rounds To At Resolution]
  \label{def:physics:phyx-mini-0907:phyxminiproblems-problemphyxmini0907-roundstoatresolution}
  \lean{PhyXMiniProblems.ProblemPhyXMini0907.RoundsToAtResolution}
  value rounds to displayed at the stated resolution when it lies strictly within half a resolution step. The last digit of 1.7 * 10⁻⁴ N has resolution 0.1 * 10⁻⁴ N = 10⁻⁵ N
\end{definition}
\begin{proof}
  This is the declared model definition of Rounds To At Resolution.
\end{proof}
\begin{definition}[Is Unique Closest Displayed Force Choice]
  \label{def:physics:phyx-mini-0907:phyxminiproblems-problemphyxmini0907-isuniqueclosestdisplayedforcechoice}
  \lean{PhyXMiniProblems.ProblemPhyXMini0907.IsUniqueClosestDisplayedForceChoice}
  \uses{def:physics:phyx-mini-0907:phyxminiproblems-problemphyxmini0907-forcemagnitudeinnewtons, def:physics:phyx-mini-0907:phyxminiproblems-problemphyxmini0907-threechargetrianglesetup, def:physics:phyx-mini-0907:phyxminiproblems-problemphyxmini0907-answerchoice, def:physics:phyx-mini-0907:phyxminiproblems-problemphyxmini0907-displayedforcemagnitudeinnewtons}
  A displayed choice is uniquely closest to the physical force magnitude
\end{definition}
\begin{proof}
  This is the declared model definition of Is Unique Closest Displayed Force Choice.
\end{proof}
\begin{lemma}[resultant force magnitude numerical bounds]
  \label{lem:physics:phyx-mini-0907:phyxminiproblems-problemphyxmini0907-resultant-force-magnitude-numerical-bounds}
  \lean{PhyXMiniProblems.ProblemPhyXMini0907.resultant_force_magnitude_numerical_bounds}
  \uses{def:physics:phyx-mini-0907:phyxminiproblems-problemphyxmini0907-forcemagnitudeinnewtons, def:physics:phyx-mini-0907:phyxminiproblems-problemphyxmini0907-threechargetrianglesetup, def:physics:phyx-mini-0907:phyxminiproblems-problemphyxmini0907-matchesproblemphysicalreadouts, def:physics:phyx-mini-0907:phyxminiproblems-problemphyxmini0907-satisfiesdisplayedtrianglegeometry, def:physics:phyx-mini-0907:phyxminiproblems-problemphyxmini0907-usesschoolcoulombconstant, def:physics:phyx-mini-0907:phyxminiproblems-problemphyxmini0907-hasphysicalelectrostaticparameters, def:physics:phyx-mini-0907:phyxminiproblems-problemphyxmini0907-calibratesdisplayeddownwardforcedirection, def:physics:phyx-mini-0907:phyxminiproblems-problemphyxmini0907-satisfiesvectorcoulombforcelaw, def:physics:phyx-mini-0907:phyxminiproblems-problemphyxmini0907-satisfieselectrostaticforcesuperposition, def:physics:phyx-mini-0907:phyxminiproblems-problemphyxmini0907-satisfiesresultantforcemagnitudelaw}
  The geometry, downward direction, Coulomb law, and superposition determine the otherwise unspecified charge q sufficiently to bound the net-force magnitude between 1.65 * 10⁻⁴ N and 1.75 * 10⁻⁴ N
\end{lemma}
\begin{proof}
  The conclusion follows from the governing relations and figure readouts named in the statement of resultant force magnitude numerical bounds.
\end{proof}
% --- Archon declaration topology end ---
% --- Archon physics formalization source end ---
